\documentclass[11pt,a4paper]{article}

\usepackage[T1]{fontenc}
\usepackage[utf8]{inputenc}
\usepackage{lmodern}
\usepackage[margin=2.5cm]{geometry}
\usepackage{microtype}
\usepackage{enumitem}
\usepackage{xcolor}
\usepackage{hyperref}

\setlength{\parindent}{0pt}
\setlength{\parskip}{0.6em}
\setlist[itemize]{topsep=0.2em,itemsep=0.25em}

\definecolor{linkblue}{HTML}{1F4E79}
\hypersetup{
    colorlinks=true,
    linkcolor=linkblue,
    urlcolor=linkblue,
    pdftitle={VIP Week 7 Report},
    pdfauthor={Yorgo Azzi}
}

\begin{document}

\begin{center}
    {\Large\bfseries VIP Week 7 Report}\\[0.5em]
    {\large Yorgo Azzi}\\[0.25em]
    October 28, 2026
\end{center}

\section{Overview}

This week, I simplified the connection between the movement-controlled music
system and Magenta RealTime 2. I also adjusted how strongly the notes produced
by the Markov system influence Magenta. The objective was to make the final
generated audio follow the movement-controlled musical output more clearly,
rather than allowing the Magenta model to behave too independently.

A smaller interface improvement was also completed: Soprano Performance mode
now uses a piano roll instead of traditional notation for the performed melody.

\section{Work Completed}

\subsection{Simplified Direct Magenta Connection}

The previous Magenta implementation used a custom communication protocol and
a separate gateway process. Although this design worked, it required two
project commands and converted the musical events into an intermediate
fixed-rate data format before converting them back to MIDI.

The performance application now sends MIDI note-on and note-off messages
directly to the virtual \texttt{Magenta RT Input} port provided by the official
Magenta application. The extra protocol and gateway files were removed. The
complete setup now requires only the Magenta standalone application and the
main hand-performance program.

The direct router still provides the two input options developed previously:

\begin{itemize}
    \item \textbf{Soprano only} sends the final melody to Magenta so that the
    model has more freedom to generate the surrounding texture.
    \item \textbf{All four voices} sends soprano, alto, tenor, and bass so that
    Magenta follows more of the harmony and rhythm created by the Markov model.
\end{itemize}

The router handles held notes, repeated pitches, note releases, pauses, and
changes between the two input options. If the Magenta port is unavailable, the
main music and camera interface continue running, while the connection is
retried in the background. An internal mock option is also retained for testing
the routing without loading the Magenta model.

\subsection{Strengthening the Influence of the Markov Output}

The complete control chain is now:

\begin{center}
Movement $\longrightarrow$ Markov music generation
$\longrightarrow$ MIDI conditioning $\longrightarrow$ Magenta audio
\end{center}

I adjusted the influence of the input notes in Magenta RealTime so that its
output responds more clearly to the pitches and timing received from the
Markov model. Because movement controls the musical decisions made by the
Markov system, this also makes the effect of the performer's movement more
recognizable in the final Magenta output.

The adjustment preserves Magenta's ability to create a richer audio texture,
while giving the input melody or four-voice texture a more noticeable role in
the result. The system therefore behaves as one connected pipeline instead of
two generators whose outputs are only loosely related.

\subsection{Soprano Piano Roll}

Soprano Performance mode now displays only the performed soprano as a live
piano roll. Pitch is shown vertically, time horizontally, and a stationary
hand produces one continuously growing note block. The complete soprano
history remains scrollable. Full SATB mode continues to use the existing
four-voice score. This visualization reads the same events sent to MIDI and
does not affect music generation or playback timing.

\section{Results}

The main results were:

\begin{itemize}
    \item Magenta now receives the selected musical voices through a direct
    MIDI connection, without a custom protocol or separate gateway process.
    \item The complete system has a simpler startup procedure and fewer
    communication components.
    \item Soprano-only and all-voice Magenta inputs remain selectable from the
    integrated interface.
    \item The influence of the input notes was adjusted so that Magenta's audio
    follows the Markov-generated pitches and timing more clearly.
    \item Movement changes can now be heard through the complete chain from
    gesture, to Markov output, to Magenta-generated audio.
    \item Soprano Performance mode now includes a compact piano-roll view,
    while Full SATB mode retains the four-voice score.
\end{itemize}

\section{Challenges}

The direct Magenta design depends on the virtual MIDI port created by the
official standalone application. The port is available only after Magenta is
running and a compatible model setup is ready. The project must also use a
different MIDI output for its immediate synthesized sound; otherwise the
ordinary performance output could be sent to the Magenta input by mistake.

Although the current pipeline works well, Magenta RealTime leaves limited room
for further technical exploration inside the generation stage. The project can
select the input notes, adjust their influence, choose between soprano and all
voices, and change the text prompt. After that, Magenta's pretrained model
performs the remaining audio generation internally. Its musical decisions and
generation process are not directly controlled by this project.

This creates a practical limitation for future development: further Magenta
work would mainly involve changing the notes and settings supplied to an
existing model, rather than designing or modifying the generation process
itself. Future work may therefore be more valuable if it focuses on the
movement mapping, the Markov model, or another accompaniment method that
offers more direct control over its musical behavior.

Although direct MIDI removes unnecessary communication steps, the perceived
response still depends on camera capture, MediaPipe tracking, the next musical
tick, MIDI synthesis, and Magenta inference. These delays must be measured
separately to determine which parts most affect the performer's experience.

\section{Next Steps}

\begin{itemize}
    \item Replace camera-based computer vision with wearable sensor-based
    movement detection. This would remove the system's dependence on lighting,
    camera position, visibility, and reliable MediaPipe tracking.
    \item Collect and stream accelerometer data from four smartwatches worn by
    one performer: one on each hand and one on each leg. The sensor signals
    would be filtered, normalized, and converted into continuous movement
    features for the realtime music model.
    \item Develop a musically logical mapping between the four limbs and the
    model's parameters. Hand movement can control melodic and expressive
    features, while leg movement can control rhythmic pulse, tempo, accents,
    or larger structural changes. The exact mapping should keep each gesture
    understandable and give the performer clear audible feedback.
\end{itemize}

\section{Resources}

\begin{itemize}
    \item Project repository:
    \url{https://github.com/YorgoAzzi/MarkovModel}
    \item Magenta RealTime 2 repository:
    \url{https://github.com/magenta/magenta-realtime}
    \item Magenta standalone application documentation:
    \url{https://github.com/magenta/magenta-realtime/blob/main/docs/apps/standalone_app.md}
    \item Qt documentation:
    \url{https://doc.qt.io/qt-5/}
\end{itemize}

\end{document}
